% Report prose: the sectioned shape. Each section of the binder is a piece;
% \pieceheading prints its title as an unnumbered heading with a contents line
% (the numbers on the headings inside it come from the body emitter, which
% never numbers a piece title, so cross-references stay in step).
%
% Restrained headings: sans-serif, bold, modest sizes. titlesec is loaded by
% preamble.tex; starred sections take the same format.
\titleformat{\section}{\sffamily\bfseries\large}{\thesection}{1em}{}
\titleformat{\subsection}{\sffamily\bfseries\normalsize}{\thesubsection}{1em}{}
\titleformat{\subsubsection}{\sffamily\itshape\normalsize}{\thesubsubsection}{1em}{}
\titlespacing*{\section}{0pt}{2.5ex plus 1ex minus .2ex}{1ex plus .2ex}
\titlespacing*{\subsection}{0pt}{2ex plus 1ex minus .2ex}{0.8ex plus .2ex}

\newcommand{\pieceheading}[1]{\section*{#1}\addcontentsline{toc}{section}{#1}}
\newcommand{\pieceheadingnotoc}[1]{\section*{#1}}

% Optional first arg `notoc` suppresses the contents line. \def, not
% \newcommand, for the comparison macro (\ifx compares the \long prefix too).
\makeatletter
\def\prose@toc{toc}
\newenvironment{prose}[2][toc]{%
  \def\prose@arg{#1}%
  \ifx\prose@arg\prose@toc \pieceheading{#2}\else \pieceheadingnotoc{#2}\fi
}{}
\makeatother

% A section break inside a piece: a short rule.
\newcommand{\scenebreak}{%
  \par\vspace{1ex}{\centering\rule{0.2\textwidth}{0.4pt}\par}\vspace{1ex}}
